% sections/simulator.tex
\section{C2: \texorpdfstring{\ffsim}{ffsim}, a run simulator}
\label{sec:sim}

The oracle answers ``what will this file score?'' for the price of a full chip run. \ffsim{} tries to answer it for
the price of a millisecond, and to say when it cannot.

\subsection{Decomposition}
\label{sec:decomp}

Because the contest charges wall-clock time, a recipe's score has two sources: how many optimizer steps fit in the
budget, and how much each step learns. \ffsim{} models them separately:
\begin{align}
  y &= q(\text{recipe}, S, \text{seed}) + \mu, \label{eq:decomp}\\
  S &= \operatorname{walk}\bigl(\text{schedule},\; T - t_{\text{res}},\; \{\bar t_k\}_{k\in\{1,2,4\}}\bigr), \notag
\end{align}
where $y$ is the official score, $q$ the rehearsal \valbpb{} after $S$ steps, $\mu$ the offset of \cref{eq:offset}
(with $\gamma_c=0$ for chip~C), $T$ the charged-time target, $t_{\text{res}}=\val{reserve-s}$\docnum{}~s the
measured stop reserve, and $\bar t_k$ the step time of the phase that accumulates $k$ micro-batches. The walk
replays the training script's clock logic: warm-up at the full batch until step~20, each accumulation phase until
its fraction of $T$, the cooldown and the stop rule. \Cref{eq:decomp} makes the trade of \cref{sec:price} explicit:
a change enters through $\bar t_k$ (and so $S$) and through $q$. The step times $\bar t_k$ come from an empirical
step-time model on code-lineage and knob features; with one code lineage left out, it predicts the $k\!=\!4$ step time
of runs whose mechanisms another lineage carried, with an MAE of \val{st-lolo-reach} (\cref{app:steptime} describes
the model and the walk).

\subsection{Data}
\label{sec:simdata}

\texttt{research/sim-data/runs.jsonl} holds \val{ds-n} run records from \val{ds-chips} Trainium instances (chips
A--D; \val{ds-nochip} records have no recorded chip): \val{ds-full} full runs and \val{ds-screens} short screens
(\cref{tab:dataset}), merged from a run table, the campaign ledger and harvested chip logs (\cref{app:simdata}).
\texttt{validation-pairs.json} holds \val{ds-pairs} treatment/control pairs from \val{ds-pair-families} families:
\val{ds-pairs-recipe} recipe changes, each at one seed on one chip, and \val{ds-pairs-noise} comparisons of identical
recipes, \val{ds-pairs-diffseed} across seeds on one chip and \val{ds-pairs-crosschip} across chips at one seed. They
were listed by hand from the campaign log of \val{pairs-listed}\docnum{} (\cref{app:simdata}). The quality model is
fitted only on the \val{q-nfit} full chip runs whose complete environment is known. The final-day runs on chips E, G
and H are not in the dataset.

\subsection{The quality surrogate}
\label{sec:quality}

The quality model is a Bayesian ridge (MAP) regression of the rehearsal score,
\begin{equation}
  \label{eq:quality}
  q = \alpha_{\text{era}} + \delta_{\text{chip}} + b L + c L^2 + \sum_j \beta_j x_j(\text{knobs}) + u_{\text{seed}} + \varepsilon_q,
  \qquad L = \ln(S/2300),
\end{equation}
where 2{,}300 steps is the typical step count of the late recipes, so the intercepts are scores at that step count.
Every coefficient has a Gaussian prior. The components are:
\begin{itemize}
  \item \textbf{Era effects} $\alpha$, one per code era (a group of code versions, \cref{app:sim}). A leakage guard
    merges an era with fewer than two runs into its parent lineage, so that a single-run treatment cannot be absorbed
    by ``its'' era.
  \item \textbf{Chip effects} $\delta$ (prior SD \val{q-prior-chip}) for every instance other than chip~C.
  \item \textbf{A step term} with a strong prior, $b\sim\mathcal{N}(\val{q-prior-b}, \val{q-prior-b-sd}^2)$ and
    $c\sim\mathcal{N}(\val{q-prior-c}, \val{q-prior-c-sd}^2)$, centred on the price of \cref{sec:price}. The prior is
    strong because within an era the step count varies mostly through mechanisms that also change quality. The
    posterior slope, \val{q-slope} per unit $L$, stays close to it and is not independent evidence for $\kappa$.
  \item \textbf{Knob features} $x_j$: \val{q-nbase-features} documented encodings of the knobs that varied, each
    centred on a reference recipe and scaled so that one unit is a typical campaign move (prior SD
    \val{q-prior-knob} bpb per unit; \cref{tab:features}), plus \val{q-nbowl-features} ``bowl'' terms that give tuned
    scalars a non-negative curvature.
  \item \textbf{Seed effects} $u$ (prior SD \val{q-prior-seed}).
\end{itemize}
The fit is a MAP solve with non-negative bowl terms and an empirical-Bayes residual SD, floored at
\val{q-sigma-floor} (\cref{app:fit}). The final $\hat\sigma_q=\val{q-sigma}$ sits on the floor; without the floor the
estimate is \val{q-sigma-nofloor}. For every candidate the model also reports
\emph{support}: whether each changed value was seen in the data, interpolated, extrapolated or never varied. In the
decision rule we used, only a supported candidate whose paired difference from the base is $-0.0003$ or better earns
a chip run. \texttt{simulate} and \texttt{search} propagate the noise by paired Monte Carlo (\cref{app:fit}).

\subsection{Validation, with baselines}
\label{sec:simval}

\Cref{tab:sim-calibration} collects the checks. Every row predicts the difference $\Delta$ between the two arms of a
pair, and sign agreement counts a predicted sign that matches the measured one. Four baselines use the same pairs and
the same leave-one-pair-out protocol: a leave-one-out majority sign, the zero predictor, the mean $\Delta$ of the
other pairs in the same family that share no run with the held-out pair (the mean of all recipe pairs for the
\val{base-fam-nfallback} pairs without one), and the same model with every knob effect switched off (era, chip, steps
and seed only). Intervals are cluster bootstraps over the \val{lopo-boot-clusters} control runs, many of which anchor
several pairs.

\begin{table}[tbp]
  \centering
  \caption{\textbf{Quality-surrogate calibration and baselines.} Sign agreement on treatment/control pairs; MAE of the
    predicted difference in $10^{-3}$ bpb. $\hat\Delta$: predicted, $\Delta$: measured difference. ``Recipe pairs''
    exclude the \val{ds-pairs-noise} seed and chip pairs, on which sign agreement is meaningless. Sign intervals are
    cluster bootstraps over control runs; for the majority sign, each pair's leave-one-out prediction is held fixed
    and the control-run clusters are resampled. The temporal rows compare the surrogate with the steps-only model and
    the zero predictor fitted and scored on the same pairs. \textsuperscript{$\dagger$}\,Per-pair values from
    \texttt{docs/simulator-report-20260929.md}, scored by \texttt{paper/analysis/simulator.py} (\cref{tab:postfit});
    its interval is a Wilson interval, too narrow because several pairs share a control.}
  \label{tab:sim-calibration}
  % GENERATED by paper/analysis -- quality surrogate calibration with baselines (simulator.py)
{\footnotesize
\begin{tabular}{@{}p{7.6cm}rrcr@{}}
\toprule
Check & Pairs & Sign & 95\% CI & MAE ($10^{-3}$) \\
\midrule
\multicolumn{5}{@{}l}{\emph{Full surrogate}} \\
\quad in-sample (reference only) & 99 & 86.9\% &  & 0.21 \\
\quad leave one pair out (LOPO), all pairs & 99 & 84.8\% & 76\%--91\% & 0.39 \\
\quad LOPO, recipe pairs only & 87 & 82.8\% & 73\%--89\% & 0.37 \\
\quad LOPO, recipe pairs with $|\hat\Delta|\ge0.3$ (decision-relevant) & 55 & 92.7\% & & \\
\quad LOPO, recipe pairs with $|\Delta|\ge0.3$ (outcome-decisive) & 52 & 96.2\% & & \\
\quad leave one knob value out (LOKVO), recipe pairs & 87 & 81.6\% &  & 0.58 \\
\multicolumn{5}{@{}l}{\emph{Baselines, same pairs and protocol}} \\
\quad steps + era + chip + seed, knob effects off (LOPO), recipe & 87 & 66.7\% & 51\%--79\% & 0.71 \\
\quad \quad recipe pairs with $|\hat\Delta|\ge0.3$ & 36 & 75.0\% & & \\
\quad leave-one-out majority sign, recipe pairs & 87 & 62.1\% & 48\%--84\% & -- \\
\quad zero predictor ($\hat\Delta=0$), recipe pairs & 87 & -- & & 0.63 \\
\multicolumn{5}{@{}l}{\emph{Forward in time and new mechanisms}} \\
\quad temporal hold-out, fit before 28 Sep & 72 & 86\% & & 0.89 \\
\quad temporal hold-out, fit before 29 Sep & 46 & 91\% & & 0.48 \\
\quad pairs finished after the release fit$^{\dagger}$ & 15 & 53\% & 30\%--75\% & 1.04 \\
\quad \quad zero predictor on the same pairs & 15 & -- & & 0.75 \\
\bottomrule
\end{tabular}
}

\end{table}

\paragraph{Leave one pair out (LOPO).} For each pair, refit without its two runs and predict their difference. On
the \val{lopo-rec-n} recipe pairs the full surrogate agrees on sign \val{lopo-rec-sign} of the time (95\% CI
\val{base-full-rec-signci}) with MAE \val{lopo-rec-mae}. The steps-only baseline scores \val{base-steps-rec-sign}
(\val{base-steps-rec-signci}) and MAE \val{base-steps-rec-mae}; the paired difference, surrogate minus baseline, is
\val{base-diff-sign-ci} of sign agreement and \val{base-diff-mae-ci}$\times10^{-3}$ of MAE. Against the family
mean the surrogate gains \val{base-diff-fam-sign-ci} of sign agreement but no clear MAE gain (\cref{app:pairval}). A majority
sign scores \val{base-major-rec-sign} and the zero predictor has MAE
\val{base-zero-rec-mae} (\cref{fig:pairs} plots every pair). Within the replicate setting of LOPO, then, the knob
features carry real information about signs. On the pairs where the \emph{prediction} is decisive, which is what
the decision rule uses, the surrogate scores \val{lopo-rec-signpdec} on \val{lopo-rec-npdec} pairs against
\val{base-steps-pdec-sign} for the steps-only model (\cref{app:pairval}). The MAE is above the noise floor of two
single runs, \val{pair-noise-floor} (\cref{app:derivations}).

\paragraph{Leave one knob value out (LOKVO).} LOPO measures replicate prediction, since the same knob values usually
stay in the fit through other seeds. LOKVO also drops every run that carries the pair's untried value, so the model
must predict a value it has never seen, which is the question the search asks. Sign agreement on recipe pairs stays
at \val{lokvo-rec-sign} (\val{lokvo-rec-signci}), but the MAE rises to \val{lokvo-rec-mae} (per family in \cref{tab:lopo-family}).

\paragraph{Forward in time.} Fit only on runs before a cutoff and predict every pair dated on or after it, as the
campaign would have (\cref{tab:temporal}). With a cutoff of 29~September the model scores \val{temp-b-sign} on
\val{temp-b-n} pairs, MAE \val{temp-b-mae}; with 28~September, fitted on \val{temp-a-nfit} runs, \val{temp-a-sign}
and MAE \val{temp-a-mae}. The steps-only model fitted on the same early runs scores \val{temp-b-steps-sign} and
\val{temp-a-steps-sign}, so the surrogate gains \val{temp-b-diff-sign-ci} and \val{temp-a-diff-sign-ci} of sign
agreement, but in MAE only with the later cutoff (\cref{app:pairval}). The temporal set excludes the incomplete-knob
runs that tested the newest mechanisms, which biases it towards supported pairs.

\paragraph{A development-set estimate.} The surrogate's specification was not frozen before these checks:
\val{sc-n}\docnum{} configurations were compared with their LOPO results visible (\cref{tab:surrogate-configs},
\cref{app:pairval}), and two feature encodings were set with campaign outcomes in view (named in
\cref{tab:features}). LOPO and LOKVO are therefore development-set estimates; the \val{postfit-n} post-fit pairs
are the only clean test.

\paragraph{Genuinely new mechanisms.} On the \val{postfit-n} pairs that finished after the release fit, which tested
attention-source variants and new learning-rate settings (\cref{tab:postfit}), sign agreement was \val{postfit-sign}
(\val{postfit-k} of \val{postfit-n}; Wilson 95\% interval \val{postfit-wilson}, too narrow because several pairs
share a control), indistinguishable from chance, and the MAE \val{postfit-mae} was worse than the zero predictor's
\val{postfit-zero-mae}. The surrogate can say which knob moves are inside the noise; it cannot say whether a new
mechanism will work. During the campaign we used it for negative screening: none of the \val{search-n}\docnum{} local
candidates around \upl{K60} cleared the chip-confirmation bar, and we ran none of them. Because they were never run,
this is not evidence that they were bad. Every later quality win came from mechanisms the data had never seen (the
row pool and the EMA blend).

\paragraph{A GPU proxy.} A second route replays the chip's optimizer trajectory on a CUDA GPU to measure per-step
quality off the chip; with $n=\val{gpu-n-inf}$ informative pairs and run-to-run noise above its agreement with the
chip, it is a feasibility demonstration, not a validated instrument (\cref{sec:gpu}).
